% Part: lambda-calculus
% Chapter: introduction
% Section: primitive-recursion

\documentclass[../../../include/open-logic-section]{subfiles}

\begin{document}

\olfileid{lam}{int}{pr}
\olsection{\usetoken{S}{lambda definable} تابعې د بنسټيز بازګښت تر عمل لاندې تړلې دي}

د بنسټيز بازګښت په برخه کښې بالاخره يو څۀ کار ته اړتيا لرو۔ دا کار
به په پړاوونو کښې کوو۔ د پخوا په شان، فرض کوو چې ترمونه $G$ او $H$
په ترتيب سره تابعې $g$ او~$h$ !!{lambda define} کوي؛ غواړو داسې
ترم $F$ ومومو چې لاندې تعريف شوې تابع~$f$ !!{lambda define}s:
\begin{align*}
f(0, \vec z) & = g(\vec z) \\
f(x+1, \vec z) & = h(x, f(x,\vec z), \vec z).
\end{align*}
\textbf{د ژباړې د نسخې سمون (OLLAM-005):} سرچينې دلته د f د
تعريفوونکي ترم نوم H راوړے و، حال دا چې H ئې لا د h دپاره
ټاکلے دے؛ د لاندې ټول ثبوت له مخې د نوي ترم نوم F دے۔
\textbf{د ژباړې د نسخې سمون (OLLAM-006):} د بازګښت د ګام په سرچينه
کښې د h لومړے آرګومېنټ z و؛ دلته د ګام شاخص x شو، چې د
راتلونکې عددنښيزې معادلې له لومړي آرګومېنټ سره هم سمون لري۔
نو په عمومي توګه، چې لامبډا ترمونه $G'$ او~$H'$ راکړل شوي وي، دومره
بس ده چې داسې ترم~$F$ ومومو چې
\begin{align*}
F(\num{0}, \vec z) & \equiv G'(\vec z) \\
F(\overline{n+1}, \vec z) & \equiv H'(\num{n}, F(\num{n}, \vec z), \vec z)
\end{align*}
د هر طبيعي عدد~$n$ دپاره پوره کړي؛ دا چې $G'$ او~$H'$ په ترتيب سره
$g$ او~$h$ !!{lambda define} کوي، معنا ئې دا ده چې کله د $\vec z$
په ځاے عددنښې $\num{\vec m}$ کښېږدو، نو $F(\num{n+1},
\num{\vec m})$ سم ځواب ته عادي کېږي۔
\textbf{د ژباړې د نسخې سمون (OLLAM-007):} سرچينې په دې دوو
معادلو کښې G او H ليکلي وو، خو دلته راکړل شوي تعريفوونکي ترمونه
G' او H' دي؛ بې‌پرېمه G او H لاندې د اضافي پارامترونو
د جذبولو دپاره بېل تعريفونه لري۔

خو د دې دپاره دومره بس ده چې داسې ترم~$F$ ومومو چې
\begin{align*}
F(\num 0) & \equiv G \\
F(\overline {n+1}) & \equiv H(\num{n},F(\num{n}))
\intertext{د هر طبيعي عدد $n$ دپاره پوره کړي، چېرته چې}
G & = \lambd[\vec z][G'(\vec z)] \text{او}\\
H(u,v) & = \lambd[\vec z][H'(u,v(\vec z),\vec z)].
\end{align*}
\textbf{د ژباړې د نسخې سمون (OLLAM-008):} په سرچينه کښې د
H(u,v) په تعريف کښې v(u,z) و؛ خو v دلته د F(n) په شان ترم دے
چې يوازې د پارامترونو لړۍ اخلي۔ له همدې امله په تعريف کښې
v(z) راغلے دے۔
په بله وينا، د لامبډا نښو په دې تدبير سره له اضافي پارامترونو
$\vec z$ سره د جلا چلند اړتيا نشته---هغوي په لامبډا نښه کښې جذبېږي۔

د ترم $F$ تر تعريفولو وړاندې د مرتبو جوړو د سمبالولو يوه وسيله پکار
ده۔ لاندې لمه دا وسيله راکوي۔

\begin{lem}
داسې لامبډا ترم $D$ شته چې د هر دوو لامبډا ترمونو $M$ او~$N$
دپاره $D(M,N)(\num{0}) \red M$ او $D(M,N)(\num{1}) \red N$ وي۔
\end{lem}

\begin{proof}
لومړے لامبډا ترم $K$ داسې تعريف کړئ:
\[
K(y) = \lambd[x][y].
\]
په بله وينا، $K$ ترم $\lambd[y][\lambd[x][y]]$ دے۔ په بل نظر، د
هر~$M$ دپاره $K(M)$ ثابته تابع ده چې پر هر ننوت~$M$ ورکوي۔

اوس $D(x,y,z)$ په $D(x,y,z) = z (K(y))x$ تعريف کړئ۔ بيا لرو
\begin{align*}
D(M,N,\num 0) & \red \num 0 (K(N)) M \red M \text{او}\\
D(M,N,\num 1) & \red \num 1 (K(N)) M \red K(N) M \red N,
\end{align*}
لکه څنګه چې غوښتل شوي وو۔
\end{proof}

مطلب دا دے چې $D(M,N)$ د مرتبې جوړې $\tuple{M, N}$ استازيتوب کوي؛
که فرض کړو چې $P$ هم د داسې جوړې استازے وي، نو $P(\num 0)$ او
$P(\num 1)$ په ترتيب سره د چپ او ښي پروجيکشنونو، $(P)_0$ او
$(P)_1$، استازيتوب کوي۔ له دې وروسته همدا وروستۍ نښې کاروو۔

\begin{lem}
!!{lambda definable} تابعې د بنسټيز بازګښت تر عمل لاندې تړلې دي۔
\end{lem}

\begin{proof}
بايد وښيو چې د هر دوو راکړل شويو ترمونو $G$ او $H$ دپاره داسې ترم
$F$ موندل کېدے شي چې
\begin{align*}
F(\num{0}) & \equiv G \\
F(\num{n+1}) & \equiv H(\num{n}, F(\num{n}))
\end{align*}
د هر طبيعي عدد~$n$ دپاره پوره کړي۔ مفکوره په لنډ ډول دا ده چې
د \emph{جوړو} لاندې لړۍ محاسبه کړو:
\[
\tuple{\num 0, F(\num 0)}, \tuple{\num 1, F(\num 1)}, \dots,
\]
او عددنښې د تکراروونکو په توګه وکاروو۔ پام وکړئ چې لومړۍ جوړه
يوازې $\tuple{\num 0, G}$ ده۔ له جوړې $\tuple{\num{n},
F(\num{n})}$ څخه راتلونکې جوړه $\tuple{\num{n+1}, F(\num{n+1})}$
بايد له $\tuple{\num{n+1}, H(\num{n}, F(\num{n}))}$ سره برابره وي۔
داسې لامبډا ترم~$T$ جوړوو چې دا يو-ګامي بدلون ترسره کړي۔

تفصيل دا دے۔ $T(u)$ داسې تعريف کړئ:
\[
T(u) = \tuple{S((u)_0), H((u)_0,(u)_1)}.
\]
اوس په آسانۍ سره کتلے شو چې د هر عدد~$n$ دپاره
\[
T(\tuple{\num{n}, M}) \red \tuple{\num{n+1}, H(\num{n}, M)}.
\]
لکه پاس چې وويل شول، د $G$ او $H$ په لرلو سره $F(u)$ داسې
تعريف کړئ:
\[
F(u) = (u(T,\tuple{\num 0, G}))_1.
\]
په بله وينا، د ننوت~$\num{n}$ دپاره $F$، $T$ پر
$\tuple{\num 0, G}$ باندې $n$ ځله تکراروي او بيا دويم غړے
ورکوي۔ په پيل کښې لرو
\begin{enumerate}
\item $\num{0} (T, \tuple{\num 0, G}) \equiv \tuple{\num{0}, G}$
\item $F(\num{0}) \equiv G$
\end{enumerate}
پر~$n$ د استقرا په وسيله د هر طبيعي عدد دپاره لاندې ښودلے شو:
\begin{enumerate}
\item $\num{n+1}(T, \tuple{\num{0}, G}) \equiv \tuple{\num{n+1}, F(\num{n+1})}$
\item $F(\num{n+1}) \equiv H(\num{n}, F(\num{n}))$
\end{enumerate}
د دويمې فقرې دپاره لرو
\begin{align*}
F(\num{n+1}) & \red (\num{n+1}(T, \tuple{\num 0, G}))_1 \\
& \equiv (T(\num{n} (T, \tuple{\num 0, G})))_1 \\
& \equiv (T(\tuple{\num{n}, F(\num{n})}))_1 \\
& \equiv (\tuple{\num{n+1}, H(\num{n}, F(\num{n}))})_1 \\
& \equiv H(\num{n}, F(\num{n})).
\end{align*}
له پای څخه په دويمه کرښه کښې مو د استقرا فرضيه کارولې ده۔ د
لومړۍ فقرې دپاره لرو
\begin{align*}
\num{n+1} (T, \tuple{\num 0, G}) &
\equiv T(\num{n} (T, \tuple{\num 0, G})) \\
& \equiv T( \tuple{\num{n}, F(\num{n})}) \\
& \equiv \tuple{\num{n+1}, H(\num{n}, F(\num{n}))} \\
& \equiv \tuple{\num{n+1}, F(\num{n+1})}.
\end{align*}
په وروستۍ کرښه کښې مو دويمه فقره کارولې ده۔ نو ومو ښودله چې
$F(\num 0) \equiv G$ او د هر $n$ دپاره $F(\num {n+1}) \equiv H(\num
n, F(\num{n}))$، او همدا زموږ غوښتنه وه۔
\end{proof}

\end{document}
